% Appendix D, first part (restructure writer W3). Proofs for sec:recovery; app-split.tex continues this appendix (no \section there).
% Must never \input app-split.tex (main.tex inputs both).
% Sources: split/iid/notes-final.md Lemmas 2.1-2.2, Thms 3.1, 3.2, 4.1, Props 3.4, 4.2-4.4, Remark 4.5, Prop 5.1, Cor 6.1, Props 7.2, 7.3, 7.5, 7.6, Lemma 8.1(iv);
% split/kink/notes-final.md Lemmas 5.1-5.3, Thm 5.4, Cor 5.7 (scheme (pq)); split/verify-iid.md §1-§2; split/verify-kink.md §2-§3.
% Local notation: Lambda_n(pi) := nu_pi(o_n)/nu(o_n) (likelihood ratio, proofs only); t_{\bar\ell} prior tail; Z^{\bar\ell}_n tail sum; g_i Hellinger gap at the i-th input
% (frozen labels: indexed by distinct inputs, DECISIONS-restructure D2). Parity classes (D11(h)): d, r, r^*, r_1, A_r, ell(r), r_0, C_{k,r_0}. rem:recovery:window stated in D.4 (RW3-05).
\section{Proofs for \texorpdfstring{\cref{sec:recovery}}{Section \ref*{sec:recovery}} (exact model recovery)}\label{app:recovery}\label{app:split}

Throughout, $\Lambda_n(\pi):=\nu_\pi(o_n)/\nu(o_n)=2^{L_n(\nu)-L_n(\pi)}$ is the likelihood ratio of $\pi$ against the sampler; almost surely every label has positive $\nu$-probability, so it is well defined. The prior's tail beyond $\bar\ell$ is $t_{\bar\ell}:=\sum_{\ell(\pi)>\bar\ell}2^{-\ell(\pi)}$, which tends to $0$, and $Z^{\bar\ell}_n:=\sum_{\ell(\pi)>\bar\ell}2^{-\ell(\pi)}\Lambda_n(\pi)$.

\subsection{Likelihood ratios}\label{app:recovery:tools}

\begin{lemma}[Likelihood ratios]\status{proved; strong law and Ville's inequality known}\src{split/iid Lemmas 2.1, 2.2(b),(c); split/kink Lemma 5.1 (pq); split/verify-kink \S2.1}\label{lem:recovery:ratios}
Let the inputs be chosen by any rule of \cref{def:split:setting}, indexed by queries for per-query labels and by distinct inputs, in order of first appearance, for frozen labels; let $o_i$ and $\Lambda_i(\pi)$ be the data and the likelihood ratio at the $i$-th step of this indexing, and put $g_i(\pi):=\Hel(\nu(\cdot\mid x_i),\nu_\pi(\cdot\mid x_i))$ for its $i$-th input, which is fixed before $y_i$ is drawn. (Indexed by queries, $\Lambda$ is constant at a repeated query with frozen labels, so the supermartingale statements of (a) hold in either indexing.)
\begin{enumerate}[label=(\alph*),itemsep=1pt,topsep=2pt]
\item $\Ex[\Lambda_{n+1}(\pi)/\Lambda_n(\pi)\mid o_n,x_{n+1}]=\sum_{y:\nu(y\mid x_{n+1})>0}\nu_\pi(y\mid x_{n+1})\le1$. So $\Lambda_n(\pi)$, $\Lambda_n(\pi)^\eta$ ($0<\eta\le1$) and $Z^{\bar\ell}_n$, also with $\Lambda^\eta$ in place of $\Lambda$, are nonnegative supermartingales, and $\Prob[\sup_nZ^{\bar\ell}_n\ge a]\le t_{\bar\ell}/a$ for $a>0$.
\item On the event $\{\sum_ig_i(\pi)=\infty\}$, almost surely $\Lambda_n(\pi)\to0$.
\end{enumerate}
Now let the data be per query with inputs i.i.d.\ from $\mu$, so that $\Lambda_n(\pi)$ is a product of $n$ independent copies of $Z:=\nu_\pi(y\mid x)/\nu(y\mid x)$, $(x,y)\sim\mu\otimes\nu$.
\begin{enumerate}[label=(\alph*),resume,itemsep=1pt,topsep=2pt]
\item $\KL_\mu(\pi)=-\Ex\log Z\in[0,\infty]$ is positive iff $\pi\notin\mathcal E$; $\Ex(\log Z)^+\le\log e/e$; and $\frac1n\log\Lambda_n(\pi)\to-\KL_\mu(\pi)$ almost surely, also when $\KL_\mu(\pi)=\infty$.
\item $\Ex\Lambda_n(\pi)^s=\mathsf C_s(\pi)^n$ for $s\in(0,1]$; for $s<1$, $\mathsf C_s(\pi)\le(1-{\Hel}_\mu(\pi))^{2\min(s,1-s)}$, so $\mathsf C_s(\pi)^n\le e^{-2\min(s,1-s){\Hel}_\mu(\pi)n}$. In particular $\Ex\sqrt{\Lambda_n(\pi)}=(1-{\Hel}_\mu(\pi))^n\le e^{-n{\Hel}_\mu(\pi)}$, and $\Lambda_n(\pi)=1$ for all $n$ and all $\pi\in\mathcal E$, almost surely.
\end{enumerate}
\end{lemma}

\begin{proof}
(a) Given $o_n$ and $x_{n+1}$, the new factor is $\nu_\pi(y\mid x_{n+1})/\nu(y\mid x_{n+1})$ with $y\sim\nu(\cdot\mid x_{n+1})$ (per query, and also with frozen labels, where the next distinct input is new and its label is drawn independently of the past). Its mean is the stated sum. Jensen's inequality gives the same for the power $\eta$, and the sums $Z^{\bar\ell}_n$ inherit the property termwise; they start at $t_{\bar\ell}$. Ville's maximal inequality \citep{ville1939etude} gives the tail bound.

(b) Given $o_n$ and $x_{n+1}$, $\Ex\sqrt{\Lambda_{n+1}/\Lambda_n}=\sum_{y:\nu>0}\sqrt{\nu(y\mid x_{n+1})\nu_\pi(y\mid x_{n+1})}=1-g_{n+1}$. If $g_{n+1}=1$, then almost surely $\nu_\pi(y_{n+1}\mid x_{n+1})=0$ and $\Lambda$ is $0$ from then on. Otherwise $M_n:=\sqrt{\Lambda_n}/\prod_{i\le n}(1-g_i)$ is a nonnegative supermartingale, so it converges almost surely, and on $\{\sum_ig_i=\infty\}$, $\prod_{i\le n}(1-g_i)\le e^{-\sum_{i\le n}g_i}\to0$, hence $\Lambda_n\to0$. This is the argument behind Kakutani's dichotomy \citep{kakutani1948equivalence}; it needs only that $g_{n+1}$ is fixed before $y_{n+1}$ is drawn.

(c) Conditional Jensen for the strictly concave logarithm gives $\Ex[\log Z\mid x]\le\log\Ex[Z\mid x]\le0$, with equality in both iff $\nu_\pi(\cdot\mid x)=c\,\nu(\cdot\mid x)$ on the support of $\nu(\cdot\mid x)$ with $c=1$, which with $\sum_y\nu_\pi\le1$ forces $\nu_\pi(\cdot\mid x)=\nu(\cdot\mid x)$. So $\KL_\mu(\pi)>0$ iff the kernels differ on a set of positive $\mu$-measure. Since $\ln z\le z/e$, $\Ex(\log Z)^+\le(\log e/e)\Ex Z\le\log e/e$, and the strong law for i.i.d.\ variables with integrable positive part gives the limit in $[-\infty,0]$.

(d) Independence gives $\Ex\Lambda_n^s=(\Ex Z^s)^n$ and $\Ex Z^s=\mathsf C_s(\pi)$. For $s\in[\frac12,1)$ write $\nu^{1-s}\nu_\pi^s=(\sqrt{\nu\nu_\pi})^{2(1-s)}\nu_\pi^{2s-1}$ and apply Hölder's inequality on $\mu\otimes$(counting measure) with exponents $1/(2(1-s))$ and $1/(2s-1)$; this gives $\mathsf C_s\le\mathsf C_{1/2}^{2(1-s)}(\Ex_\mu\sum_y\nu_\pi)^{2s-1}\le\mathsf C_{1/2}^{2(1-s)}$. For $s\in(0,\frac12]$ write $\nu^{1-s}\nu_\pi^s=(\sqrt{\nu\nu_\pi})^{2s}\nu^{1-2s}$ to get $\mathsf C_s\le\mathsf C_{1/2}^{2s}$. Then use $\mathsf C_{1/2}=1-{\Hel}_\mu$ and $1-u\le e^{-u}$. For $\pi\in\mathcal E$ every $x_i$ falls almost surely where the two kernels agree, and there are countably many programs.
\end{proof}

\subsection{Proof of \texorpdfstring{\cref{thm:recovery:main}}{Theorem \ref*{thm:recovery:main}}}\label{app:recovery:main}

Work on the almost sure event of \cref{lem:recovery:ratios}(d) on which $\Lambda_n(\pi)=1$ for all $n$ and all $\pi\in\mathcal E$. A program $\pi_0\in\mathcal E_0$ then has $J_\lambda(\pi_0;h_n)=\lambda\ell_0+L_n(\nu)<\infty$, so $\lambda$-MAPs exist for every $\lambda>0$ (\cref{lem:setup:basic}(i)).

\emph{(a).} If $\hat\pi$ is a $\lambda_n$-MAP, then $J_{\lambda_n}(\hat\pi)\le J_{\lambda_n}(\pi_0)$, that is,
\begin{equation}\label{eq:recovery:beats}
\log\Lambda_n(\hat\pi)\ \ge\ \lambda_n\big(\ell(\hat\pi)-\ell_0\big).
\end{equation}
\emph{Long programs, uniformly in $n$.} Let $\bar\ell\ge\ell_0$. If a $\lambda_n$-MAP has $\ell(\hat\pi)>\bar\ell$, then \eqref{eq:recovery:beats} and $\lambda_n\ge1$ give $\log\Lambda_n(\hat\pi)\ge\ell(\hat\pi)-\ell_0$, so $2^{-\ell(\hat\pi)}\Lambda_n(\hat\pi)\ge2^{-\ell_0}$ and $Z^{\bar\ell}_n\ge2^{-\ell_0}$. By \cref{lem:recovery:ratios}(a), $B_{\bar\ell}:=\{\exists n:Z^{\bar\ell}_n\ge2^{-\ell_0}\}$ has probability at most $2^{\ell_0}t_{\bar\ell}\to0$. Since $Z^{\bar\ell}_n$ is pathwise nonincreasing in $\bar\ell$, the events $B_{\bar\ell}$ decrease, and their intersection is null. So almost surely some finite $\bar\ell_0$ bounds the length of every $\lambda_n$-MAP at every $n$.
\emph{Short programs.} By \cref{lem:recovery:ratios}(c) and countability, almost surely $\frac1n\log\Lambda_n(\pi)\to-\KL_\mu(\pi)<0$ for all $\pi\notin\mathcal E$ at once. Take $\pi\notin\mathcal E$ with $\ell(\pi)\le\bar\ell_0$. If $\ell(\pi)\ge\ell_0$, the left side of \eqref{eq:recovery:beats} tends to $-\infty$ and the right side is nonnegative. If $\ell(\pi)<\ell_0$, the left side is eventually at most $-n\KL_\mu(\pi)/2$ (or tends to $-\infty$ if $\KL_\mu(\pi)=\infty$), while the right side is $-\lambda_n(\ell_0-\ell(\pi))=-o(n)$. Programs of $\mathcal E\setminus\mathcal E_0$ have $\Lambda_n=1$ and $\ell>\ell_0$, so they violate \eqref{eq:recovery:beats} at every $n$. Finitely many programs have length at most $\bar\ell_0$, so from some $n$ on every $\lambda_n$-MAP lies in $\mathcal E_0$, and all of $\mathcal E_0$ tie.

\emph{(b).} Divide the posterior weights by $\nu(h_n)$; they become $2^{-\lambda\ell(\pi)}\Lambda_n(\pi)$, and the normaliser is at least $W_\lambda$. For every $\bar\ell$, since $2^{-\lambda\ell}\le2^{-\ell}$,
\[
\sum_{\pi\notin\mathcal E}2^{-\lambda\ell(\pi)}\Lambda_n(\pi)\ \le\ \sum_{\pi\notin\mathcal E,\ \ell(\pi)\le\bar\ell}2^{-\lambda\ell(\pi)}\Lambda_n(\pi)+Z^{\bar\ell}_n .
\]
The finite sum tends to $0$ almost surely, by \cref{lem:recovery:ratios}(c). Choose $\bar\ell_j$ with $t_{\bar\ell_j}\le4^{-j}$; then $\Prob[\sup_nZ^{\bar\ell_j}_n\ge2^{-j}]\le2^{-j}$, and by the Borel--Cantelli lemma almost surely $\sup_nZ^{\bar\ell_j}_n<2^{-j}$ for all large $j$. So the mass outside $\mathcal E$ tends to $0$, the normaliser tends to $W_\lambda$, and the terms of $\mathcal E$ are constant. For the tempered posterior use the weights $2^{-\ell(\pi)}\Lambda_n(\pi)^\eta$, $\Lambda_n^\eta\to0$ off $\mathcal E$, and the tail bound for $\Lambda^\eta$; the normaliser tends to $W_1$.

\emph{(c).} A program with $J_\lambda(\pi)\le J_\lambda(\pi_0)$ has $\lambda(\ell(\pi)-\ell_0)\le\log\Lambda_n(\pi)\le\ell(\pi)+\delta_n$ by \cref{lem:split:ville}, which is the length bound. A fill-in has $\Lambda_n(\pi)=2^{L_n(\nu)}\le2^{\ell(\pi)+\delta^*}$, which is its length bound. Hence, summing over fill-ins,
\[
\sum2^{-\lambda\ell(\pi)}\Lambda_n(\pi)\le2^{-(\lambda-1)(L_n(\nu)-\delta^*)}\sum2^{-\ell(\pi)}\Lambda_n(\pi)\le2^{-(\lambda-1)(L_n(\nu)-\delta^*)+\delta^*},
\]
and the normaliser divided by $\nu(h_n)$ is at least $2^{-\lambda\ell_0}$. For $\Pi^{1/\lambda}$ a fill-in has $\Lambda_n^{1/\lambda}=2^{-(1-1/\lambda)L_n(\nu)}\Lambda_n$, so the fill-ins' weight is at most $2^{-(1-1/\lambda)L_n(\nu)+\delta^*}$, against a normaliser of at least $2^{-\ell_0}$. Finally $L_n(\nu)/n\to H_\mu(\nu)$ by the strong law, since $0\le\Ex[-\log\nu(y\mid x)]\le\log|Y|$.

\emph{(d).} Almost surely the programs of $\mathcal E\setminus\mathcal E_0$ are worse than $\pi_0$ by $\lambda(\ell-\ell_0)>0$, so the event that the set of $\lambda$-MAPs is not $\mathcal E_0$ is contained in $\bigcup_{\pi\notin\mathcal E}\{\Lambda_n(\pi)\ge2^{\lambda(\ell(\pi)-\ell_0)}\}$. For each such $\pi$ and each $s\in(0,1]$, Markov's inequality for $\Lambda_n(\pi)^s$ and \cref{lem:recovery:ratios}(d) bound its probability by $2^{-s\lambda(\ell(\pi)-\ell_0)}\mathsf C_s(\pi)^n$; take the infimum over $s$ and a union bound. For the margin form split the union at $\bar\ell:=\bar\ell_\varepsilon$. A program with $\ell(\pi)>\bar\ell$ and $J_\lambda(\pi)\le J_\lambda(\pi_0)$ forces $\delta^*>(\lambda-1)\bar\ell-\lambda\ell_0\ge\log(2/\varepsilon)$ by (c), which has probability at most $\varepsilon/2$. For $\ell(\pi)\le\bar\ell$ take $s=1/\lambda$: since $\min(1/\lambda,1-1/\lambda)=\min(1,\lambda-1)/\lambda$, the sum is at most $\sum_{\ell(\pi)\le\bar\ell}2^{\ell_0-\ell(\pi)}e^{-an}\le2^{\ell_0}e^{-an}$ with $a:=2\min(1,\lambda-1)\gamma/\lambda$, and this is at most $\varepsilon/2$ when $an\ge(\ell_0+\log(2/\varepsilon))\ln2$. \qed

\subsection{The parity classes, and what is necessary}\label{app:recovery:lambda1}

\emph{Parities.} In $\mathcal C_k$ and $\mathcal C_{k,r_0}$ a parity on $r$ bits is consistent with $h_n$ iff the label vector $y$ lies in the column span over GF$(2)$ of the $n\times r$ matrix $A_r$ of the first $r$ bits of the inputs, and a parity on fewer bits is one on more bits. The labels are fair coins independent of the inputs, so $\Prob[y\in\mathrm{colspan}A_r]=\Ex2^{\mathrm{rank}A_r-n}\le2^{r-n}$. With i.i.d.\ uniform rows, $\Prob[\mathrm{rank}A_r<n]\le\sum_{i<n}2^{i-r}<2^{n-r}$. The coin has $J_\lambda=\lambda k+n$, and a consistent parity has $J_\lambda$ equal to $\lambda$ times its length, which is $\ell(r):=r+2\lceil\log(r+1)\rceil+1$ on $r$ bits in $\mathcal C_k$; inconsistent parities have $J_\lambda=\infty$. A parity has $\mathsf C_s=2^{s-1}$.

\emph{Proof of \cref{prop:recovery:lambda1}(a).} At $n=r_0-5$, $\Prob[\mathrm{rank}A_{r_0}<n]<2^{-5}$, and when the rank is $n$ some parity is consistent, with $J_1=r_0+2=n+7<k+n$, so every $1$-MAP is a parity.

\emph{(b), lower end.} For $n$ with $2\lceil\log(n+6)\rceil+6<k$ and $r=n+5\le d$, $\ell(r)=n+6+2\lceil\log(n+6)\rceil<k+n$, and with probability at least $1-2^{-5}$ a parity on $r$ bits is consistent and strictly beats the coin. At $n_b:=2^{(k-8)/2}-6$ ($k$ even), $2\lceil\log(n_b+6)\rceil+6=k-2<k$, so the probability is at least $31/32>\varepsilon$ at $n_b$, and $n_{\rm rec}(\varepsilon)>n_b$.

\emph{(b), upper end.} Let $r^*:=\max\{r\ge1:\ell(r)\le k+n\}$; the set of $1$-MAPs differs from $\{\text{coin}\}$ (ties included) iff $y\in\mathrm{colspan}A_{r^*}$, so its probability $P(n)$ is at most $2^{r^*-n}$ (capping $r^*$ at $d$ only lowers it). From $\ell(r^*)\le k+n$, $r^*-n\le k-1-2\lceil\log(r^*+1)\rceil$, so $P(n)\le2^{k-1}/(r^*+1)^2$. Put $r_1:=n+k-1-2\lceil\log(n+1)\rceil$. If $1\le r_1\le n$, then $\ell(r_1)\le n+k$ and $r^*\ge r_1$; if $r_1>n$, then $\ell(n)<n+k$ and $r^*\ge n$. For $n\ge32$, $n/2\ge2\lceil\log(n+1)\rceil$, so $r_1\ge n/2$ and $r^*+1\ge n/2$. Hence $P(n)\le2^{k+1}/n^2$ for $n\ge32$, a decreasing bound that is at most $\varepsilon$ once $n\ge2^{(k+1)/2}/\sqrt\varepsilon$.

\emph{(c).} \Cref{thm:recovery:main}(d) with $s=1/\lambda$ for every parity gives at most $2^k2^{(1/\lambda-1)n}$ times the total prior weight of the parities, which is at most $1/4$ in either class.

The ingredients of \cref{thm:recovery:main}(d) are necessary in the following sense.

\begin{proposition}[What is necessary]\status{proved}\src{split/iid Props 4.2, 4.3, 4.4(a),(b), Remark 4.5; split/verify-iid \S1.5}\label{prop:recovery:necessity}
(a) \emph{A margin.} Let $\lambda>1$, $\ell_0\ge2$ and $\lambda(\ell_0-1)\ge9$. For every $N$ some class of two programs, one of them a program of length $\ell_0$ for $\nu$, has $\Prob[\text{the set of $\lambda$-MAPs at $n$}\ne\mathcal E_0]\ge1-8.66/(\lambda(\ell_0-1))$ for every $n\le N$.
(b) \emph{The factor $1/\gamma$.} Let $\nu$ be the fair coin with a program of length $\ell_0$, and $\pi'$ a program for $\Bern(\frac12+\varepsilon')$, $0<\varepsilon'\le\frac1{10}$, of length $\ell_0-\Delta\ge1$, with $\lambda\Delta\ge9$. Then ${\Hel}_\mu(\pi')\in[\varepsilon'^2/2,2\varepsilon'^2]$, and in the class of these two programs the $\lambda$-MAP is $\pi'$ with probability at least $1-8.66/(\lambda\Delta)$ for every $n\le0.083\,\lambda\Delta/{\Hel}_\mu(\pi')$.
(c) \emph{The factor $1/(\lambda-1)$.} In $\mathcal C_k$ with $\lambda>1$, for every $n$ with $(\lambda-1)n<\lambda(k-6-2\lceil\log(n+6)\rceil)$ and $d\ge n+5$, with probability at least $31/32$ the $\lambda$-MAP is a parity and the coin is not a $\lambda$-MAP.
(d) \emph{The length bound is attained.} Let $k$ be an integer and $\lambda>1$ with $n:=\lambda(k-2)/(\lambda-1)$ a positive integer. In $\mathcal C_{k,n}$, with probability $\prod_{i=1}^n(1-2^{-i})>0.2887$ the set of $\lambda$-MAPs at $n$ is $\{\text{coin},\varphi\}$ for a parity $\varphi$ with $\ell(\varphi)=(\lambda k+\delta_n)/(\lambda-1)-\log(1+2^{2-k})/(\lambda-1)$.
(e) \emph{The canonical class.} Let $Y=\{0,1\}$, $\mathcal C=\Hcl{\times}$, and $\nu$ computable with $\nu(1\mid x)\in[\beta,1-\beta]$ for all $x$ and a rational $\beta\in(0,\frac12]$. For every computable $g:\mathbb N\to(0,1]$, $\gamma_{\bar\ell}<g(\bar\ell)$ for all large $\bar\ell$.
\end{proposition}

\emph{Proof of \cref{prop:recovery:necessity}(b).} The $\lambda$-MAP is $\pi'$ iff $\log\Lambda_n(\pi')>-\lambda\Delta$, and $\log\Lambda_n(\pi')=\sum_{i\le n}X_i$ with $X_i\in\{\log(1+2\varepsilon'),\log(1-2\varepsilon')\}$ equally likely. The mean is $\Ex X_1=\frac12\log(1-4\varepsilon'^2)\ge-2\varepsilon'^2/((1-4\varepsilon'^2)\ln2)$ and $\Var X_i\le\log^2(1-2\varepsilon')\le4\varepsilon'^2/((1-2\varepsilon')^2\ln^22)$. For $n\le n_1:=\lambda\Delta(1-4\varepsilon'^2)\ln2/(4\varepsilon'^2)$ we have $n|\Ex X_1|\le\lambda\Delta/2$, and Chebyshev's inequality gives
\[
\Prob\Big[\sum_iX_i\le-\lambda\Delta\Big]\le\frac{4n\Var X_1}{(\lambda\Delta)^2}\le\frac{4(1+2\varepsilon')}{(1-2\varepsilon')\ln2\cdot\lambda\Delta}\le\frac{8.66}{\lambda\Delta}.
\]
Since $\sqrt{1+u}+\sqrt{1-u}\le2-u^2/4$, ${\Hel}_\mu(\pi')\ge\varepsilon'^2/2$; and ${\Hel}_\mu(\pi')\le\frac12\sum_y(P-Q)^2/P=2\varepsilon'^2$. So $n_1\ge0.96\ln2\cdot\lambda\Delta/(8{\Hel}_\mu(\pi'))\ge0.083\lambda\Delta/{\Hel}_\mu(\pi')$. \emph{(a)} is (b) with $\Delta=\ell_0-1$ and $\varepsilon'$ so small that $0.083\lambda(\ell_0-1)/{\Hel}_\mu(\pi')\ge N$; the Kraft sum is $\frac12+2^{-\ell_0}\le1$.

\emph{(c).} With $r=n+5$, $\lambda\ell(r)<\lambda k+n$ is the stated condition, and the rank argument applies.

\emph{(d).} If the $n\times n$ matrix $A_n$ is invertible, which has probability $\prod_{i=1}^n(1-2^{-i})$, exactly one parity $\varphi$ is consistent. Then $J_\lambda(\varphi)=\lambda(n+2)$ equals $J_\lambda(\text{coin})=\lambda k+n$ by the choice of $n$, and all other parities have $J_\lambda=\infty$. The mixture is $\xi_{\mathcal C}(h_n)=2^{-k}2^{-n}+2^{-(n+2)}$, so $\delta_n=\log(2^{-k}+\frac14)=-2+\log(1+2^{2-k})$, and $\ell(\varphi)=n+2=(\lambda k-2)/(\lambda-1)$ is the stated value.

\emph{(e).} Let $\pi_0\in\mathcal E_0$, and let $\tau$ be the flipped kernel, $\tau(y\mid x):=\nu_{\pi_0}(1-y\mid x)$, or $\tau:=\Bern(\frac14)$ if $\nu$ is the fair coin $\mu$-a.e.; either way $\tau$ is lsc and differs from $\nu$ on a set of positive $\mu$-measure (which case holds need not be decided). The kernels $\nu_j:=(1-2^{-j})\nu_{\pi_0}+2^{-j}\tau$ are lsc uniformly in $j$ and outside $\mathcal E$, and by closure under prefix-free simulation they have programs of length at most $\ell_0+K(j)+\kappa$. Since $\Hel(P,Q)\le\frac12\sum_y(P-Q)^2/P$, $|\nu_j-\nu|\le2^{-j}$ and $\nu\ge\beta$, their gaps are at most $2^{-2j}/\beta$. Approximating $g(\bar\ell)$ until an approximation $a\ge3\cdot2^{-i}$ appears gives a computable rational $g'(\bar\ell):=a-2^{-i}\in[g(\bar\ell)/2,g(\bar\ell)]$, and $h(m):=\lceil\frac12\log(2/(\beta g'(m+\ell_0+\kappa)))\rceil+1$ is computable with $2^{-2h(m)}/\beta<g(m+\ell_0+\kappa)$. Let $j(\bar\ell)$ be the largest $j$ with $K(j)\le\bar\ell-\ell_0-\kappa$; since $K(h(m))\le2\log(m+1)+\kappa'\le m$ for large $m$, $j(\bar\ell)\ge h(\bar\ell-\ell_0-\kappa)$, and $\gamma_{\bar\ell}\le2^{-2j(\bar\ell)}/\beta<g(\bar\ell)$ for all large $\bar\ell$. \qed

\subsection{Proof of \texorpdfstring{\cref{thm:recovery:kink}}{the sample-kink theorem} and of \texorpdfstring{\cref{rem:recovery:window}}{the window remark}}\label{app:recovery:kink}

\begin{remark}[The $\ell_0$ in the right window]\status{proved; a natural strengthening refuted}\src{split/iid Lemma 8.1(iv); split/verify-iid \S1.7}\label{rem:recovery:window}
The natural strengthening of \cref{thm:recovery:kink}(iii) that each bit of length buys at most one bit of fit plus $\delta^*$ in total, $\rho(\ell_0)-\rho(\alpha)\le(\alpha-\ell_0)+\delta^*$ for the sample's structure function $\rho$, fails: in $\mathcal C_{k,r_0}$ with $r_0+2>k$ and $n=r_0-5\ge1$, with probability at least $31/32$ a parity is consistent and $\rho(\ell_0)-\rho(r_0+2)-(r_0+2-\ell_0)=k-7$, so it fails with probability at least $31/32-2^{7-k}$. Hence no fixed $\lambda\le1$ lies in a kink interval guaranteed at finite $n$; only the limit in \cref{thm:recovery:kink}(v) puts $\lambda=1$ inside.
\end{remark}

\emph{Proof of \cref{thm:recovery:kink}(i).} Fix $\pi$ and put $a(\pi):=2\log e\cdot n{\Hel}_\mu(\pi)-2\ell(\pi)-2\log(1/\varepsilon)$. By Markov's inequality and \cref{lem:recovery:ratios}(d), $\Prob[\Lambda_n(\pi)\ge2^{-a(\pi)}]=\Prob[\sqrt{\Lambda_n(\pi)}\ge2^{-a(\pi)/2}]\le2^{a(\pi)/2}e^{-n{\Hel}_\mu(\pi)}=\varepsilon2^{-\ell(\pi)}$. A union bound over $\pi$ costs at most $\varepsilon\sum_\pi2^{-\ell(\pi)}\le\varepsilon$, and on the complement $L_n(\pi)-L_n(\nu)=-\log\Lambda_n(\pi)>a(\pi)$.

\emph{(ii).} Programs shorter than $\ell_0$ are outside $\mathcal E$, so on the event of (i) they have $L_n(\pi)-L_n(\nu)>2\log e\cdot n\gamma_=(\ell(\pi))-2\ell(\pi)-2\log(1/\varepsilon)\ge\sigma_n(\ell_0-\ell(\pi))$, and the programs of $\mathcal E_0$ have fit $L_n(\nu)$, so they bound the point. For the second form, $(A-2\ell)/(\ell_0-\ell)\ge A/\ell_0$ iff $\ell(A-2\ell_0)\ge0$.

\emph{(iii).} By \cref{lem:recovery:ratios}(a), $\Ex[\xi_{\mathcal C}(h_n)/\nu(h_n)]=\sum_\pi2^{-\ell(\pi)}\Ex\Lambda_n(\pi)\le1$, so by Markov's inequality, with probability at least $1-\varepsilon$, $2^{-\ell(\pi)}\Lambda_n(\pi)<1/\varepsilon$ for every $\pi$, that is, $L_n(\pi)>L_n(\nu)-(\ell(\pi)-\ell_0)-(\ell_0+\log(1/\varepsilon))$. This is at least $L_n(\nu)-\tau(\ell(\pi)-\ell_0)$ once $(\tau-1)(\ell(\pi)-\ell_0)\ge\ell_0+\log(1/\varepsilon)$, for every $\tau>1$ on the same event. The bound for all $n$ is \cref{lem:split:ville}; the offset $\ell_0+\delta^*$ is nonnegative because $\delta^*\ge\delta_n\ge-\ell_0$.

\emph{(iv).} On the event of (i), a program outside $\mathcal E$ of length at most $\bar\ell$ has $L_n(\pi)-L_n(\nu)>2\log e\cdot n\gamma_{\bar\ell}-2\bar\ell-2\log(1/\varepsilon)>0$ by the hypothesis on $n$; programs of $\mathcal E$ have fit $L_n(\nu)$. This is the flat part, and every right secant to a length in $(\ell_0,\bar\ell]$ is at most $0$. Beyond the cap, on the event of (iii), a secant is below $(\ell+\log(1/\varepsilon))/(\ell-\ell_0)\le\tau_{\bar\ell}$ for $\ell>\bar\ell$. With (ii), $(\ell_0,L_n(\nu))$ is a zero-window $(\sigma_n,\tau_{\bar\ell})$-kink with $\tau_{\bar\ell}<\sigma_n$, and \cref{cor:kink:locate}(c) makes it a vertex whose programs are the $\lambda$-MAPs for $\lambda\in(\tau_{\bar\ell},\sigma_n)$; by the flat part these programs form $\mathcal E_0$. The probability is at least $1-2\varepsilon$.

\emph{(v).} Apply (i) with $\varepsilon_n:=n^{-2}$; by the Borel--Cantelli lemma its event holds for all large $n$, and then $s_-(n)\ge\min_{\ell<\ell_0}(2\log e\cdot n\gamma_=(\ell)-2\ell-4\log n)/(\ell_0-\ell)$; divide by $n$. By \cref{thm:recovery:main}(a) with $\lambda_n\equiv1$, almost surely the set of $1$-MAPs is eventually $\mathcal E_0$, whose programs all sit at the point $(\ell_0,L_n(\nu))$. By \cref{lem:setup:hull}(v), a set of $\lambda$-MAPs at a single point is the set of programs at a vertex $v_i$ with $\sigma_{i+1}<\lambda<\sigma_i$, since at $\lambda=\sigma_i$ both ends of an edge are minimisers. The largest right secant from $v_i$, floored at $0$, is $\sigma_{i+1}$ (it is $0$ at a last vertex), and by convexity the left steepness at $v_i$ is at least $\sigma_i$. So $s_+(n)=\sigma_{i+1}<1<\sigma_i\le s_-(n)$.

\emph{\Cref{rem:recovery:window}.} In $\mathcal C_{k,r_0}$ with $r_0+2>k$ the coin is the only program of length at most $k=\ell_0$, so $\rho(\ell_0)=L_n(\text{coin})=n$; a consistent parity has fit $0$ and length $r_0+2$, so the excess is $n-(r_0+2-k)=k-7$. Some parity is consistent with probability at least $\Prob[\mathrm{rank}A_{r_0}=n]\ge1-2^{n-r_0}=31/32$, and $\Prob[\delta^*\ge k-7]\le2^{7-k}$. \qed

\subsection{Proof of \texorpdfstring{\cref{cor:recovery:noisy}}{the corollary on noisy computable functions}}\label{app:recovery:noisy}

(i) $\pi\in\mathcal E$ iff $\nu_\pi(\cdot\mid x)=\nu_{f,p}(\cdot\mid x)$, for both labels, at $\mu$-a.e.\ $x$. For $p<\frac12$ the labels have probabilities $1-p>p$, so such a kernel determines $f(x)$ and $p$; apply \cref{thm:recovery:main}(a). (ii) is \cref{thm:recovery:main}(c), with $H(\nu_{f,p}(\cdot\mid x))=H(p)$ at every $x$. (iii) For $p=\frac12$, $\nu_{f,1/2}$ is the coin for every $f$. For $p=0$, $\nu_{f,0}$ is the point mass on $f$, the programs that compute $f$ $\mu$-a.e.\ are in $\mathcal E$, and $L_n(\nu)=0$. Finally $\nu_{f,p}(1\mid x)\in\{p,1-p\}\subseteq[\beta,1-\beta]$ for every rational $\beta\in(0,p]$. \qed

\begin{remark}[Frozen labels on a finite domain]\status{proved}\src{split/iid Cor 6.1(vi); split/verify-iid \S1.10}\label{rem:recovery:frozentable}
Let $\mu$ have finite support $X_0$, with frozen labels. The data set is eventually the full table $F$ on $X_0$, so every $\lambda$-MAP is eventually a minimiser of $J_\lambda(\cdot\,;F)$, a function of $F$ alone, and by \cref{prop:recovery:frozen}(a) no procedure identifies every such sampler. Among the programs $(f,q)$ with the true $f$, the one of least $J_\lambda$ has the rate $q$ that minimises $\lambda\ell(f,q)+K\log(1/q)+(|X_0|-K)\log(1/(1-q))$, $K$ the number of flipped inputs (a penalised empirical rate); the $\lambda$-MAP is this program unless another function or the coin has smaller $J_\lambda$. In an explicit class (inputs uniform on $\{0,1\}^6$, parities $a$ with rates $q$ from an eight-element list with a prefix code, and a coin; the truth $(a^*,\frac18)$ has length $10$ and $(a^*,0)$ length $9$), on the event that no input is flipped, of probability $(7/8)^{64}$, the program $(a^*,0)$ has $J_\lambda=9\lambda<10\lambda+64\log(8/7)=J_\lambda(a^*,\tfrac18)$; so for every $\lambda$ the $\lambda$-MAP is eventually wrong with probability at least $(7/8)^{64}$.
\end{remark}

A direct computation (code in the accompanying repository) gives the exact failure probabilities $0.255$, $0.255$, $0.204$ and $0.566$ in this class at $\lambda=1$, $1.5$, $2$ and $4$, and $0.136$, $0.136$, $0.085$ and $0.566$ in a second class, with another rate code and single-input exception programs; per query \cref{cor:recovery:noisy}(i) applies to both classes.

\subsection{Frozen labels, adaptive inputs, misspecification and \texorpdfstring{$\lambda=cn$}{lambda=cn}}\label{app:recovery:beyond}

\emph{The run with frozen labels.} With frozen labels and i.i.d.\ inputs on the countable $X$, a program of $\mathcal E$ equals $\nu$ at every input of the support of $\mu$, which contains every queried input almost surely, so it can serve as $\pi_\nu$ in \cref{thm:split:private}. If the support is infinite, the distinct inputs eventually enumerate it, and since the terms are nonnegative, $\nu$-equivalence on the run is $\sum_{x\in\supp\mu}\Hel(\nu(\cdot\mid x),\nu_\pi(\cdot\mid x))<\infty$. Per query with i.i.d.\ inputs, $g_i(\pi)$ are i.i.d.\ in $[0,1]$ with mean ${\Hel}_\mu(\pi)$, so for $\pi\notin\mathcal E$ their sum is infinite almost surely by the strong law, and for $\pi\in\mathcal E$ it is $0$.

\emph{Proof of \cref{prop:recovery:frozen}(a).} With frozen labels the data process is a measurable function of the input sequence and of the table $F$ on the support of $\mu$, drawn once with independent values $F(x)\sim\nu(\cdot\mid x)$. The input sequence has the same law under both samplers and is independent of $F$. The laws of $F$ under $\nu$ and $\nu'$ are products that differ in the single factor at $x_0$, with density $\nu'(F(x_0)\mid x_0)/\nu(F(x_0)\mid x_0)\in(0,\infty)$ almost surely; so the laws of the data process are mutually absolutely continuous. Since $\mu(x_0)>0$, $\mathcal E_\nu\cap\mathcal E_{\nu'}=\emptyset$. If $A:=\{\hat\pi_n\in\mathcal E_\nu\text{ eventually}\}$ has probability $1$ under $\nu$, it has probability $1$ under $\nu'$, and on $A$ the estimator is eventually outside $\mathcal E_{\nu'}$. In $\Hcl{\times}$ take $\nu'(\cdot\mid x_0):=\frac12(\nu(\cdot\mid x_0)+\ind[\cdot=y_0])$, which is computable, distinct from $\nu(\cdot\mid x_0)$ and has the same support.

\emph{(b).} A deterministic $\varphi$ has $\Hel(\nu(\cdot\mid x),\ind[\cdot=\varphi(x)])=1-\sqrt{\nu(\varphi(x)\mid x)}\ge1-\sqrt{1-\beta_x}\ge\beta_x/2$, and gap $1$ where it is undefined. So no deterministic program is $\nu$-equivalent on the run, and \cref{thm:split:private}(iv) applies.

\emph{(c).} The table is eventually complete. An $e_j$ with $F(j)=0$ has likelihood $0$. An $e_j$ with $F(j)=1$ fits $\log64=6$ bits better than $\nu$ and is $4$ bits longer, so it beats $\nu$ iff $4\lambda<6$; all such $e_j$ tie. Some $F(j)$ is $1$ with probability $1-(63/64)^{10}$, and otherwise $\nu$'s program is the $\lambda$-MAP.

For the limits of \cref{thm:recovery:main} in directions other than frozen labels, put $\KL_\mu(\pi):=\Ex_{x\sim\mu}\KL(\nu(\cdot\mid x)\,\Vert\,\nu_\pi(\cdot\mid x))\in[0,\infty]$, the mean KL divergence, and $R_\nu(\pi):=H_\mu(\nu)+\KL_\mu(\pi)=\Ex_{\mu\otimes\nu}[-\log\nu_\pi(y\mid x)]$, the expected per-query fit.

\begin{proposition}[Adaptive inputs, misspecification, $\lambda=cn$]\status{proved}\src{split/iid Props 7.2(b), 7.3(iii$'$), 3.4(a)--(c); split/verify-iid \S1.3, \S1.8}\label{prop:recovery:beyond}
\begin{enumerate}[label=(\alph*),itemsep=1pt,topsep=2pt]
\item \emph{Adaptive inputs.} Let the inputs be $x_j=0^j$, let $\nu(1\mid0^j)=\frac12+2^{-j/2-2}$ have a program of length $k$, and let the class consist of it and the coin, of length $k-\Delta\ge1$ with $\Delta>0$. The coin is $\nu$-equivalent on the run although it differs from $\nu$ at every query, and for every fixed $\lambda>0$, with positive probability, it is the $\lambda$-MAP for all large $n$.
\item \emph{Misspecification: no fill-ins per query.} Let $\nu$ be a computable $f$ with independent flips at rate $p\in(0,\frac12)$, per query with i.i.d.\ inputs, and let the class consist of deterministic programs and one fair coin, so that $\mathcal E=\emptyset$. The coin, the unique minimiser of $\KL_\mu$, is almost surely eventually the unique $\lambda$-MAP for every $\lambda>0$.
\item \emph{$\lambda=cn$.} Per query, in a class as in \cref{def:split:setting}, let $J_c(\pi):=c\,\ell(\pi)+R_\nu(\pi)$ and suppose $J_c(\pi)<\infty$ for some $\pi$. Then $A_c:=\arg\min J_c$ is nonempty and finite, and almost surely every $cn$-MAP is in $A_c$ for all large $n$; if the programs of $A_c$ have one kernel $\mu$-a.e., the set of $cn$-MAPs equals $A_c$ for all large $n$. If $\mathcal E\ne\emptyset$ and $c^*:=\min\{\KL_\mu(\pi)/(\ell_0-\ell(\pi)):\ell(\pi)<\ell_0\}$ ($c^*:=\infty$ if no program is shorter than $\ell_0$), then almost surely the set of $cn$-MAPs is eventually $\mathcal E_0$ for $c<c^*$, and eventually contains no program of $\mathcal E_0$ for $c>c^*$; $c=c^*$ is not covered.
\end{enumerate}
\end{proposition}

With adaptive inputs the target is $\nu$ on the run, and the selected program is still a sampler (\cref{thm:split:private}(iii)); with $\lambda=cn$ selection follows the population hull of the points $(\ell(\pi),R_\nu(\pi))$ (\cref{lem:setup:hull}).

\emph{Proof of \cref{prop:recovery:beyond}(a).} The factors $\nu_{\rm coin}(y_j)/\nu(y_j\mid0^j)$ of $\Lambda_n(\text{coin})$ are independent and positive with mean $1$, and $\sum_j\Hel(\nu(\cdot\mid0^j),\text{coin})\le\sum_j2\cdot2^{-j-4}<\infty$. By Kakutani's theorem \citep{kakutani1948equivalence} the product converges almost surely and in $L^1$ to a limit $\Lambda_\infty>0$ with $\Ex\Lambda_\infty=1$, so $\Prob[\Lambda_\infty\ge1]>0$. On $\{\Lambda_\infty>2^{-\lambda\Delta}\}$, eventually $\log\Lambda_n(\text{coin})>-\lambda\Delta$, that is, $J_\lambda(\text{coin})<J_\lambda(\nu)$. The Kraft sum is $2^{-k}+2^{-(k-\Delta)}\le1$.

\emph{(b).} The coin has $\KL_\mu=1-H(p)<\infty$, and every deterministic program has $\KL_\mu=\infty$ since $p>0$. An atom of positive mass is queried infinitely often, and its labels are i.i.d.\ with both values of probability at least $p$, so almost surely some input eventually receives both labels; from then on every deterministic program has likelihood $0$, and the coin, with finite $J_\lambda$, is the unique $\lambda$-MAP.

\emph{(c).} $J_c(\pi)\ge c\ell(\pi)$, so only finitely many programs have $J_c$ below any finite value, and $A_c$ is nonempty and finite. Put $\hat R_n(\pi):=L_n(\pi)/n$, so $J_{cn}(\pi;h_n)/n=c\ell(\pi)+\hat R_n(\pi)$; by the strong law for nonnegative i.i.d.\ variables, possibly of infinite mean, $\hat R_n(\pi)\to R_\nu(\pi)$ almost surely for all $\pi$ at once. Fix $\pi_c\in A_c$. A $cn$-MAP has $cn\,\ell(\hat\pi)\le J_{cn}(\pi_c)$, so $\ell(\hat\pi)\le\ell(\pi_c)+\sup_n\hat R_n(\pi_c)/c$, which is finite almost surely because $\hat R_n(\pi_c)$ converges to the finite $R_\nu(\pi_c)$. Each of the finitely many candidates outside $A_c$ has $J_{cn}/n\to J_c>\min J_c$, so it is eventually not a MAP. Programs of $A_c$ with one kernel $\mu$-a.e.\ have equal $R_\nu$, hence equal lengths, and almost surely equal likelihoods, so they tie. If $\mathcal E\ne\emptyset$, then for $\pi_0\in\mathcal E_0$, $J_c(\pi)-J_c(\pi_0)=c(\ell(\pi)-\ell_0)+\KL_\mu(\pi)$: on $\mathcal E$ it is nonnegative with equality iff $\pi\in\mathcal E_0$; off $\mathcal E$ with $\ell(\pi)\ge\ell_0$ it is positive; and for $\ell(\pi)<\ell_0$ it is $\KL_\mu(\pi)-c(\ell_0-\ell(\pi))$, positive if $c<c^*$. So $A_c=\mathcal E_0$ for $c<c^*$. If $c>c^*$, some $\pi'$ with $\ell(\pi')<\ell_0$ has $\frac1n[J_{cn}(\pi_0)-J_{cn}(\pi')]\to c(\ell_0-\ell(\pi'))-\KL_\mu(\pi')>0$. \qed
